\talkslide{attack-objectives}{Defining the jailbreak objective}{Timing: provisional}{GCG and AdvPrompter optimize target-response likelihood. AdvGame judges sampled responses.}{
\begin{tikzpicture}[remember picture,overlay,
 every node/.style={text=Ink,align=left,inner sep=0pt,font=\normalsize}]
\begin{scope}[shift={(current page.south west)}]
 \node[anchor=north west,font=\normalsize\bfseries] at (.65,7.05) {Target strings};
 \node[anchor=north west,text=Muted] at (.65,6.5) {(e.g., GCG and AdvPrompter)};
 \node[anchor=north west,text width=6.7cm] at (.65,5.6)
  {Increase the likelihood of a chosen \textbf{affirmative target response}};
 \node[anchor=north west,font=\normalsize\itshape] at (.65,4.3)
  {``Sure, here is\ldots''};
 \draw[Muted!35,line width=.6pt] (7.95,3.2)--(7.95,7.05);
 \node[anchor=north west,font=\normalsize\bfseries] at (8.6,7.05) {Judge-based feedback};
 \node[anchor=north west,text=Muted] at (8.6,6.5) {(e.g., Self-RedTeam, AdvGame)};
 \node[anchor=north west,font=\small] at (8.6,5.6)
  {\safetyjudgetable{6.7cm}};
\end{scope}
\end{tikzpicture}


}
\speak{GCG and AdvPrompter optimize a target-response likelihood objective. The target can be an affirmative opening rather than a full harmful answer. A judge can instead assess the response the defender actually produces. The table recalls Beam's safety criterion as a concrete example of behavioral judgment, not the literal rubric of every method listed. We will distinguish the two players' compliance and deflection objectives later. PAIR already used judge feedback in 2023, so this is a contrast between objectives, not a claim that judges are new. AdvPrompter also uses judges in some evaluations, which is different from its training objective. In AdvGame, judge feedback trains both players. Faithfulness, compliance versus deflection, and the judge's own failure modes then matter for both updates. Sources: GCG, Section 2.1; AdvPrompter, Sections 2.1, 3.1, and 4.3; \href{https://arxiv.org/abs/2310.08419}{PAIR}, Section 3; AdvGame, Section 3. The next examples examine how a judge scores a response. Compliance and deflection are different criteria; after establishing how to compare responses, we will choose which criterion each player should optimize. Source detail: \href{https://arxiv.org/abs/2307.15043}{GCG}, Section 2.1; \href{https://arxiv.org/abs/2404.16873}{AdvPrompter}, Sections 2.1 and 3.1; \href{https://proceedings.mlr.press/v306/paulus26b.html}{AdvGame}, Section 3.}
